% Conditioning-only control floor: KRR-RBF((G, D, Y)) -> observable(impact) fit on
% train, evaluated on test_b and test_c. This is the parametric floor any learned
% encoder must clear to demonstrate that its latent (not the parameters alone) is
% doing the work.
% Source: outputs/session18/exp_b1_test3/conditioning_only_baseline.csv.

\begin{table}[t]
  \centering
  \caption{Conditioning-only control floor. A kernel-ridge regressor with an
  RBF kernel is fit on the train split mapping the gust parameters
  $\mathbf{c} = (G, D, Y)$ directly to each physical observable at the impact frame,
  with no latent representation. The resulting $R^2$ is the floor that any learned
  encoder $+$ predictor must clear on held-out cases to demonstrate that the latent
  representation is contributing beyond what is already implicit in the parameters.
  The rightmost column reports the JEPA d $= 64$ train $R^2$ from Table~\ref{tab:b1_closure_train_r2}
  for direct comparison.}
  \label{tab:conditioning_floor}
  \small
  \begin{tabular}{l c c c c c}
    \toprule
    Observable & $n_{\text{train}}$ & train $R^2$ & test\_b $R^2$ & test\_c $R^2$ & JEPA d=64 train $R^2$ \\
    \midrule
    $C_L$            & 226 & 0.895 & 0.303  & 0.547  & 0.864 \\
    $C_D$            & 226 & 0.941 & 0.386  & -2.459 & 0.802 \\
    $I_y$            & 226 & 0.422 & -0.301 & -1.335 & 0.562 \\
    wake-enstrophy   & 226 & 0.806 & 0.482  & -1.252 & 0.934 \\
    circ.\ pos.\     & 226 & 0.792 & -0.032 & -0.249 & 0.922 \\
    circ.\ neg.\     & 226 & 0.799 & 0.568  & -4.262 & 0.927 \\
    \bottomrule
  \end{tabular}
\end{table}
